% sec-11: Document 11, emergency preparedness and business continuity.
\docpart{Emergency Plan}{PDAC LLM Oncology Clinical Trial Site Emergency
Preparedness and Business Continuity Plan}

\section{Scope and Command Structure}

\draftnote{Name the incident commander, the alternate, and the succession, using
the roles fixed in document 14 \sect{}1 rather than inventing titles. State the
one decision that may not be delegated, which is the decision to stop a
procedure in progress.}

\section{Fire and Evacuation}

\draftnote{Alarm, response, evacuation routes, assembly, and the special case of
a participant under anesthesia. Cite NFPA 101 and keep the routes identical to
the egress plan in document 08 \sect{}1.}

\section{Seismic and Structural Events}

\draftnote{Immediate actions, the post-event inspection sequence, and the
condition on which robotic operations may resume. State that resumption requires
a re-verification of the robot registration, and cite the registration
requirement in document 06 \sect{}4.}

\section{Power Failure and Utility Loss}

\begin{mvtable}
\begin{tabularx}{\textwidth}{>{\raggedright\arraybackslash}p{3.0cm} >{\raggedright\arraybackslash}p{2.4cm} >{\raggedright\arraybackslash}p{2.4cm} >{\raggedright\arraybackslash}X}
\headrow \thead{System} & \thead{Backup} & \thead{Runtime} & \thead{Action at exhaustion}\\
\midrule
Procedure suite & to be fixed & to be fixed & to be fixed \\
Machine room & to be fixed & to be fixed & to be fixed \\
\end{tabularx}
\end{mvtable}
\tabcapl{tab:power}{Backup power by system. Stage 2 fixes the runtimes against
the essential electrical provisions in document 08 \sect{}3, and states the
action at exhaustion for every row.}

\section{Cybersecurity Incident and Model Fault}

\draftnote{Two distinct events with two distinct responses: a network compromise
and a model producing output outside its verified envelope. Take the fault
taxonomy from paper 11 of the deck table and the response structure from the NIST
Cybersecurity Framework 2.0. State that a model fault does not stop the trial; it
stops the advisory output and the human pathway continues, which is the whole
point of the advisory-only boundary in document 07 \sect{}5.}

\section{Business Continuity, Drills, and Training}

\draftnote{Recovery point and recovery time objectives, the mutual aid
arrangement for participant transfer, the quarterly drill rotation, and the
training requirement. Keep the training hours identical to document 12
\sect{}5.}
